\subsection{The first surviving seven-element source opposite a right-zero corner}

\begin{theorem}[Proved: a triangular seven-element source opposite $R_3$]
\label{CP2:A-seven-source-three-right-zero}
Let $P=\{1,e_{00},e_{01},e_{10},e_{11}\}$ be the identity-adjoined
$2\times2$ rectangular band, put $e=e_{00}$, and let
$M=P\amalg\{a,b\}$ have, in addition to the band products,
\[
 ak=bk=ek,\qquad ka=kb=ke\quad(k\in P\setminus\{1\}),
 \qquad a^2=ab=e,\quad ba=a,\quad b^2=b.
\]
Let $N=R_3$ be the identity-adjoined two-element right-zero band.
For every finite nonempty left-$N$, right-$M$ biset $S$, bounded
rational homology of its collage forces all rational degree-two
products to vanish.  There is no restriction on $|S|$.

The monoid $M$ is source $3$ in the frozen ordinary positive
order-seven catalog used by the eighteen-arrow scope.  Thus its
entire source contribution to $(7,3,8)$ is excluded by this theorem.
\end{theorem}
\begin{proof}
We construct an actual quotient collage and compute the possible
homological difference from its bar map.  Define a unital monoid
retraction
\[
 \pi:M\longrightarrow P,\qquad \pi(a)=\pi(b)=e,
 \qquad \pi|_P=\operatorname{id}_P.
\]
Every displayed product respects $\pi$.

Write $A(s)=sa$, $B(s)=sb$ and $E(s)=se$, with function
composition read from right to left.  Then
\[
 B^2=B,\quad E^2=E,\quad A^2=E,\quad AB=A,\quad BA=E,
 \quad AE=EA=BE=EB=E.
\]
In particular $ES\subseteq AS\cap BS$, and conversely a state
$t\in AS\cap BS$ satisfies $t=Bt=Et$; hence
$AS\cap BS=ES$.  All three images are invariant under $N$.

First form the $N$-equivariant quotient set $T$ by identifying
$t$ with $Et$ for every $t\in BS$.  Its canonical representatives
are $(S\setminus BS)\amalg ES$.  It has a right $P$ action:
for a band element $k$ and $t=Bs$,
$tk=sbk=sek=(Et)k$.  Next form $U$ by identifying, in $T$,
each state $u\in AS$ with $Eu$.  Its representatives are
\[
 U=(S\setminus(AS\cup BS))\amalg ES.
\]
The additional identifications are $N$ invariant and preserved
by $P$, since $ak=ek$ for every band element $k$.
The quotient map $q_S:S\to U$ satisfies
$q_S(As)=q_S(Bs)=q_S(Es)$, so is compatible with $\pi$ on
all of $M$.  Thus it defines an actual functor
\[
 q:C=(M,N,S)\longrightarrow C'=(P,N,U)
\]
which fixes the two objects and their target corner $N$.

We next compute the actual bridge comparison.  Put
$R=\mathbb QM$, $R_0=\mathbb QP$ and
\[
 h=b-e,\qquad z=a-e,\qquad p=1-h.
\]
Direct multiplication gives
\[
 h^2=h,\quad z^2=0,\quad hz=z,\quad zh=0,
 \qquad h\mathbb QK=\mathbb QKh=z\mathbb QK=\mathbb QKz=0,
\]
where $K=P\setminus\{1\}$.  Moreover
\[
 Rh=\mathbb Qh,\qquad
 Rp=\mathbb Q\{p,e_{00},e_{01},e_{10},e_{11},z\},
 \qquad pRp\xrightarrow[\pi]{\cong}R_0.
\]
The last corner has identity $p$; this is an algebra identification,
not an assertion that $p$ is a monoid element.  Let
$\Delta=R_0/\mathbb QK$, a one-dimensional right $R_0$ module.
There is an exact sequence of left-$R$, right-$R_0$ bimodules
\[
 0\longrightarrow Rh\otimes_{\mathbb Q}\Delta
   \xrightarrow{\ h\otimes1\mapsto z\ }Rp
   \xrightarrow{\pi}R_0\longrightarrow0.                 \tag{*}
\]
Both terms on the left are projective as left $R$ modules.
The right $R_0$ action on $Rp$ is through $pRp$; the displayed
map is right-linear because $zK=0$ and $zp=z$.
The stated bases prove exactness.

Let $L=\mathbb Q[S]$, with its commuting left $N$ action.
Tensoring (*) with $L$ represents $L\otimes_R^{\mathbf L}R_0$
by the natural right-$R_0$ complex
\[
 F_L=[\,Lh\otimes_{\mathbb Q}\Delta
                    \xrightarrow{\ v\mapsto vz\ }Lp\,]
\]
in degrees one and zero.  Derived associativity then identifies
the right-action bar of $S$ over $M$ with
$F_L\otimes_{R_0}^{\mathbf L}\mathbb Q$.
These are natural comparisons of the actual maps induced by
$\pi$ and $q_S$, and are left-$\mathbb QN$ linear throughout.

The projection $Lp\to\mathbb Q[T]$ is an integral, hence rational,
module isomorphism with basis
\[
 [s]-[Bs]+[Es]\quad(s\notin BS),\qquad [c]\quad(c\in ES).
\]
Indeed the projection of $t\in BS$ is $[Et]$, and the other
basis vectors have separate outside-$BS$ coordinates.  Likewise
$Lh$ has basis $[t]-[Et]$ for $t\in BS\setminus ES$.
Under these identifications its differential sends this vector to
\[
 [At]-[Et]\in\mathbb Q[T].
\]
Because $A(BS)=AS$ and $AS\cap BS=ES$, the image is exactly the
span of $[u]-[Eu]$ for $u\in AS\setminus ES$.  These vectors
are linearly independent by their separate outside-$BS$
coordinates.  The cokernel is therefore the actual state module
$\mathbb Q[U]$.  Write $W$ for the kernel on the degree-one
coefficient space $Lh$.  It is a finite-dimensional left
$\mathbb QN$ module, since all maps commute with $N$.
The kernel of the natural chain surjection
$F_L\to\mathbb Q[U]$ is consequently quasi-isomorphic to
$(W\otimes_{\mathbb Q}\Delta)[1]$.

For clarity, the second shift can be computed from the band
resolution rather than an amplitude estimate.  The explicit
length-two projective resolution in
Lemma~\ref{CP2:A-band-spheres} has the terms
$R_0e$, $R_0(1-e)$, $R_0e$ in degrees zero, one, two,
respectively, for this $2\times2$ band.  Tensoring with
$\Delta$ kills both $R_0e$ terms and leaves one copy of
$\mathbb Q$ in degree one.  Thus
\[
 \Delta\otimes_{R_0}^{\mathbf L}\mathbb Q\simeq\mathbb Q[1].
\]
It follows that the homotopy fiber in rational chains of the
actual bridge map
\[
 B(N,S,M)\longrightarrow B(N,U,P)
\]
is $(\mathbb Q\otimes_{\mathbb QN}^{\mathbf L}W)[2]$.
In particular its mapping cone has homology
$\operatorname{Tor}^{\mathbb QN}_{j-3}(\mathbb Q,W)$ in
degree $j$, with negative indices interpreted as zero.
Taking $L=\mathbb Q$ in (*) gives $Lh=0$, so the corner map
$BM\to BP$ is an integral homology equivalence as well;
the same resolution argument may be used over $\mathbb Z$.

Here are the consequences for the collage, including the shift
and the low-degree transfer.  In the two-block chain
mapping-cylinder model, bridge chains enter one degree above
the two vertex complexes.  The mapping cone of the map of
vertex complexes is acyclic, by the preceding corner comparison
and the identity on $BN$.  The mapping-cone long exact sequence
for the two collage models therefore identifies the relative
group in degree $j$ with the bridge mapping-cone group in
degree $j-1$.  The actual functor $q$ consequently has exact
sequence
\[
 \cdots\longrightarrow H_j(BC;\mathbb Q)
 \longrightarrow H_j(BC';\mathbb Q)
 \longrightarrow\operatorname{Tor}^{\mathbb QN}_{j-4}(\mathbb Q,W)
 \longrightarrow H_{j-1}(BC;\mathbb Q)\longrightarrow\cdots .
                                                               \tag{**}
\]
This exhibits both shifts from the state kernel and both shifts
from the bridge and collage cones.  In particular $q$ induces
an isomorphism on $H_2$, hence a natural isomorphism on rational
$H^2$.

Finally, the trivial right module of $N=R_3$ has projective
dimension at most one.  If its band elements are $u,v$, put
$f=1-u$.  The exact right-projective resolution is
\[
 0\longrightarrow f\mathbb QN
   \xrightarrow{\ f\mapsto v-u\ }u\mathbb QN
   \xrightarrow{\varepsilon}\mathbb Q\longrightarrow0:
\]
the first module has basis $f$, the second basis $u,v$,
and its augmentation kernel is the span of $v-u$.
Thus the Tor term in (**) vanishes except possibly in
degrees $j=4,5$.  Boundedness of $BC$ implies boundedness
of $BC'$.  Corollary~\ref{CP2:A-baseline-right-zero-three-six}
applies to $C'$, whose other corner is $P_5$, and kills all
its rational degree-two products.  Every second class of
$BC$ pulls back from $BC'$ by the established $H^2$
isomorphism, so its products vanish too.

The literal catalog identification is recorded in
\path{converse/scope18/seven_source_three_literal_binding.json}:
identity $0$, band $\{1,2,3,4\}$ with $e=1$, and $a=5,b=6$
give every product in source row $3$ of the frozen catalog.
This proves the stated source-scope consequence.
\end{proof}
